\section{Conclusiones}
La predicción de la siguiente palabra permite estudiar de forma concreta el procesamiento de secuencias. El contexto aporta las palabras anteriores y el modelo asigna probabilidades a posibles continuaciones. La salida depende tanto de la arquitectura como de la representación del texto, el vocabulario y los ejemplos utilizados durante entrenamiento.

La RNN simple mantiene un estado oculto que se actualiza en cada posición. La LSTM añade una celda y compuertas que regulan la retención, la incorporación y la salida de información. Esta diferencia explica su mecanismo de memoria, pero no garantiza una mejora automática en cualquier experimento. En la implementación, ambas arquitecturas reinician sus estados por ventana y solo acceden a los elementos incluidos en ella.

Se construyó un procedimiento que utiliza documentos separados para entrenamiento, validación y prueba, un vocabulario común de 3000 entradas y 100000 objetivos seleccionados para aprender. El programa implementa las dos redes, una referencia de bigramas, métricas y una demo. Las pruebas de integridad verifican propiedades de preparación, evaluación y recuperación del modelo necesarias para interpretar la ejecución.

Sobre 56796 objetivos de prueba, la RNN obtuvo perplejidad media de 75.60 y la LSTM de 78.15, frente a 85.93 del bigrama. El top-5 léxico fue 26.58\%, 26.14\% y 25.62\%, respectivamente. La comparación favorece a la RNN en las condiciones fijadas y muestra una mejora limitada del acierto en las primeras cinco posiciones. No demuestra superioridad general de RNN sobre LSTM.

La tasa de desconocidas de 21.19\% restringe la capacidad de identificar palabras concretas. Los ejemplos de la demo muestran esa limitación con términos del entorno local y del área de software. El sistema es útil para estudiar recurrencia y evaluación, pero los resultados no respaldan presentarlo como un autocompletador de amplia cobertura listo para un uso productivo.

La perplejidad y el acierto deben interpretarse conjuntamente. La primera resume probabilidad asignada a los objetivos, incluida la categoría desconocida; el segundo exige identificar una palabra conocida en una posición determinada. Diferenciar ambas medidas evita confundir una mejora de probabilidad con un aumento equivalente de palabras acertadas.

Como extensión del mismo problema, se podría estudiar un vocabulario de mayor cobertura, un presupuesto de entrenamiento diferente o más repeticiones de la comparación de contextos. Esas propuestas requieren nuevos protocolos y evaluaciones independientes. No forman parte de los resultados realizados ni se utilizan para sustituir las limitaciones del experimento actual.

\section{Enlaces a recursos}
El código del proyecto está organizado en scripts de Python y los registros permiten consultar las métricas sin redondear. La guía de aprendizaje complementaria explica las rutas, los comandos y la procedencia de cada elemento. Las direcciones públicas del repositorio y de la presentación se incorporarán cuando los recursos estén publicados y se haya comprobado su acceso. No se presenta una dirección no creada como si estuviera disponible.

El corpus utilizado puede consultarse en
\url{https://github.com/UniversalDependencies/UD_Spanish-AnCora/tree/r2.17}.
Los enlaces gratuitos a las doce publicaciones se encuentran en las referencias y en la guía complementaria. La presentación queda fuera de la revisión actual por indicación del grupo.
